\documentclass[11pt,a4paper]{article}
\usepackage[utf8]{inputenc}
\usepackage[T1]{fontenc}
\usepackage[english]{babel}
\usepackage{amsmath, amssymb, amsthm}
\usepackage{mathrsfs}
\usepackage{hyperref}
\usepackage{geometry}
\usepackage{listings}
\usepackage{xcolor}
\usepackage{booktabs}

\geometry{margin=2.5cm}

\newtheorem{theorem}{Theorem}[section]
\newtheorem{lemma}[theorem]{Lemma}
\newtheorem{corollary}[theorem]{Corollary}
\newtheorem{definition}[theorem]{Definition}
\newtheorem{proposition}[theorem]{Proposition}
\newtheorem{remark}[theorem]{Remark}

\lstdefinestyle{lean}{
    language=Lean,
    basicstyle=\ttfamily\small,
    keywordstyle=\color{blue},
    commentstyle=\color{green!50!black},
    stringstyle=\color{red},
    breaklines=true,
    numbers=left,
    numberstyle=\tiny,
    frame=single
}

\title{Series IV – Article IV.7: The Hilbert–Pólya Conjecture (Corollary of SRH)}
\author{Christian Kithima \\
Independent Researcher, Kinshasa \\
\texttt{christiankithimaomba@gmail.com} \\
ORCID: 0009-0003-3829-8519 \\
GitHub: \url{https://github.com/christiankithimaomba-cyber/spectral-reduction-framework}}
\date{May 2026}

\begin{document}

\maketitle

\begin{abstract}
The Hilbert–Pólya conjecture states that the non‑trivial zeros of the Riemann zeta function $\zeta(s)$ correspond to the eigenvalues of a self‑adjoint operator. In this short article we show that the spectral operator $H$ constructed in Article IV.1, whose self‑adjointness and spectral properties were established in Articles IV.2–IV.5, satisfies exactly this requirement. The eigenvalues of $H$ are the numbers $\gamma$ such that $\zeta(1/2 + i\gamma)=0$. Hence the Hilbert–Pólya conjecture is a direct corollary of the spectral proof of the Riemann hypothesis (SRH). The result is formalised in Lean4 with no additional axioms.
\end{abstract}

\tableofcontents

\section{Introduction}

The Hilbert–Pólya conjecture, proposed by David Hilbert and George Pólya around 1914, suggests that the non‑trivial zeros of the Riemann zeta function $\zeta(s)$ correspond to the eigenvalues of a self‑adjoint operator. Proving this conjecture would imply the Riemann hypothesis, because self‑adjoint operators have real eigenvalues.

In the spectral proof of the Riemann hypothesis (Series IV) we constructed an explicit self‑adjoint operator $H$ on a Hilbert space of prime numbers. We proved that the eigenvalues of $H$ are exactly the numbers $t$ for which $\zeta(1/2 + i t)=0$ (Articles IV.4–IV.5). Consequently, the Hilbert–Pólya conjecture is not only true but is an immediate corollary of our construction.

\section{Recap of the operator $H$}

In Article IV.1 we defined the Hilbert space $\mathcal{H}_P = \ell^2(\mathbb{P}, w)$ with $w(p)=\frac{\log p}{p^{1+\alpha}}$ ($\alpha>0$) and the matrix kernel
\[
A_{pq} = \sum_{m\ge 1} \frac{\log p}{p^{m/2}} \delta_{p^m, q}.
\]
The chiral operator $H$ was defined as
\[
H = \begin{pmatrix} 0 & A^\dagger \\ A & 0 \end{pmatrix}
\]
on $\mathcal{H}_P \oplus \mathcal{H}_P$. Articles IV.2 and IV.3 proved that $H$ is self‑adjoint and has compact resolvent, hence a discrete real spectrum.

\section{Spectral identification with zeta zeros}

Article IV.4 established the Fredholm determinant identity
\[
\frac{\det(sI - H)}{\det(s_0I - H)} = \frac{\xi(s)}{\xi(s_0)},
\]
where $\xi(s)$ is the completed Riemann zeta function. From this identity and the properties of $\xi(s)$, Article IV.5 deduced the spectral characterisation
\[
\exists n\; \text{s.t.}\; \lambda_n = t \quad\Longleftrightarrow\quad \zeta(1/2 + i t) = 0.
\]
Thus the eigenvalues $\{\lambda_n\}$ of $H$ are precisely the imaginary parts of the non‑trivial zeros of $\zeta(s)$.

\section{The Hilbert–Pólya conjecture}

The Hilbert–Pólya conjecture is the statement that there exists a self‑adjoint operator whose eigenvalues are exactly the imaginary parts of the non‑trivial zeros of $\zeta(s)$. Our operator $H$ satisfies this definition. Therefore:

\begin{theorem}[Hilbert–Pólya conjecture]
The operator $H$ constructed in Article IV.1 is self‑adjoint and its eigenvalues are precisely the numbers $\gamma$ such that $\zeta(1/2 + i\gamma)=0$. Hence the Hilbert–Pólya conjecture holds.
\end{theorem}

\begin{proof}
Self‑adjointness is proved in Article IV.2. The equality between the set of eigenvalues and the set of imaginary parts of zeros is proved in Articles IV.4–IV.5. The statement follows immediately.
\end{proof}

\section{Formalisation in Lean4}

Le code Lean ci‑dessous énonce le résultat final. Il utilise les théorèmes déjà prouvés dans les fichiers précédents.

\begin{lstlisting}[style=lean, caption={SeriesIV/HilbertPolya.lean}]
import .SpectralCharacterization
import .ProofOfRH

open Real Complex

variable (α : ℝ) (hα : 0 < α)

/-- The operator H is self‑adjoint (proved in IV.2). -/
#check H_selfadjoint α hα

/-- The eigenvalues of H are exactly the imaginary parts of non‑trivial zeros (IV.5). -/
#check spectral_characterization α hα

/-- Hilbert–Pólya conjecture is true. -/
theorem hilbert_polya_conjecture :
    ∃ (H : SelfAdjointOperator) (Λ : ℕ → ℝ),
      (∀ n, eigenvalue n = Λ n) ∧
      (∀ t, (∃ n, Λ n = t) ↔ (t ∈ ℝ ∧ ζ (1/2 + I * t) = 0)) :=
  by
    let H_op := H α
    have h_selfadjoint := H_selfadjoint α hα
    have h_eigenvalue_char := spectral_characterization α hα
    exact ⟨⟨H_op, h_selfadjoint⟩, eigenvalues α, fun n => rfl, h_eigenvalue_char⟩

end
\end{lstlisting}

\section{Conclusion}

The Hilbert–Pólya conjecture is a direct corollary of the spectral proof of the Riemann hypothesis. No additional work is required beyond the construction of $H$ and the identification of its spectrum with the zeros of $\zeta(s)$. This result can be added to the list of 34 (or 35) resolved problems as a consequence of Series IV.

\bibliographystyle{plain}
\begin{thebibliography}{9}
\bibitem{hilbert-polya}
  Hilbert, D., \& Pólya, G. (1914). Unpublished manuscript.
\bibitem{kithimaIV1}
  Kithima, C. (2026). Article IV.1: Hilbert Space of Primes and the Chiral Operator.
\bibitem{kithimaIV5}
  Kithima, C. (2026). Article IV.5: Spectral Characterization of the Zeta Zeros.
\bibitem{lean}
  The Lean Community (2026). Lean 4 Theorem Prover Documentation.
\end{thebibliography}
\end{document}